% Part: first-order-logic
% Chapter: syntax-and-semantics
% Section: introduction

\documentclass[../../../include/open-logic-section]{subfiles}

\begin{document}

% SEGMENT OLP-0645-S01
\olfileid{fol}{syn}{int}

\olsection{அறிமுகம்}

% SEGMENT OLP-0645-S02
முதல்தர தருக்கத்தின் கோட்பாட்டையும் மேல்கோட்பாட்டையும் வளர்க்க, அதன்
வெளிப்பாடுகளின் தொடரமைப்பையும் பொருண்மையியலையும் முதலில் வரையறுக்க வேண்டும்.
முதல்தர தருக்கத்தின் வெளிப்பாடுகள் சொற்களும் !!{formula}s உம் ஆகும். சொற்கள்
!!{variable}s, !!{constant}s மற்றும் !!{function}s ஆகியவற்றிலிருந்து
உருவாக்கப்படுகின்றன. !!^{formula}s, மறுபுறம், !!{predicate}s மற்றும் சொற்களுடன்
சேர்ந்து உருவாகின்றன (இவை மிகச் சிறிய, ``அணு'' !!{formula}s); பின்னர் அணு
!!{formula}s இலிருந்து தருக்கத் தொடுப்புக்குறிகளையும் அளவையடைகளையும் கொண்டு
மிகச் சிக்கலானவற்றை உருவாக்கலாம். உருவாக்க விதிகளை எழுதுவதற்குப் பல வழிகள்
உள்ளன; அவற்றில் ஒரு சாத்தியமான வழியை மட்டுமே தருகிறோம். பிற அமைப்புகள்
வேறு குறிகளைத் தேர்ந்தெடுக்கலாம், முதன்மையான தொடுப்புக்குறிகளின் வேறு
தொகுப்பைத் தேர்ந்தெடுக்கலாம், அடைப்புக்குறிகளை வேறுவிதமாகப் பயன்படுத்தலாம்
(அல்லது போலிஷ் குறியீட்டில் போல அவற்றைப் பயன்படுத்தாமலும் இருக்கலாம்).
எல்லா அணுகுமுறைகளுக்கும் பொதுவானது என்னவென்றால், உருவாக்க விதிகள் சொற்களின்
மற்றும் !!{formula}s இன் தொகுப்பை \emph{தொகுத்தறிதல் வழியாக} வரையறுக்கின்றன.
அது முறையாகச் செய்யப்பட்டால், ஒவ்வொரு வெளிப்பாடும் உருவாக்க விதிகளின்படி
அடிப்படையில் ஒரே ஒரு வழியில் மட்டுமே பெறப்படும். வெளிப்பாடுகள்
\emph{தனித்த வாசிப்புடையவை} ஆகும் என்ற இந்தத் தொகுத்தறிதல் வரையறை,
அதே முறையை---தொகுத்தறிதல் வரையறையைப் பயன்படுத்தி---அவற்றுக்கு
பொருள்களையும் வழங்க வழி செய்கிறது.

% SEGMENT OLP-0645-S03
வெளிப்பாடுகளுக்குப் பொருள் வழங்குவது பொருண்மையியலின் களம். பொருண்மையியலின்
மையக் கருத்து !!a{structure} இல் நிறைவுறுத்தல். !!^a{structure} மொழியின்
கட்டுமானக் கூறுகளுக்குப் பொருள் வழங்குகிறது: !!a{domain} என்பது வெறுமையற்ற
பொருட்களின் தொகுப்பு. அளவையடைகள் இப்பொருட்களத்தில் இயங்குவதாகப்
பொருள்படுத்தப்படுகின்றன; !!{constant}s பொருட்களத்தின் உறுப்புகளுக்கு
ஒதுக்கப்படுகின்றன; !!{function}s !!{domain} இலிருந்து அதற்கே செல்லும்
சார்புகளுக்கு ஒதுக்கப்படுகின்றன; !!{predicate}s !!{domain} மீது உள்ள
தொடர்புகளுக்கு ஒதுக்கப்படுகின்றன. !!{domain} மற்றும் அடிப்படைச்
சொற்களஞ்சியத்திற்கான ஒதுக்கீடுகள் சேர்ந்து !!a{structure} ஐ உருவாக்குகின்றன.
!!^{variable}s !!{formula}s இல் தோன்றலாம்; அவற்றுக்கு பொருண்மையியல்
வழங்க, !!{domain} இன் !!{element}s ஐ அவற்றுக்கும் ஒதுக்க வேண்டும்---இதுவே
மாறி ஒதுக்கீடு. இறுதியாக, நிறைவுறுத்தல் தொடர்பு இக்கூறுகளை ஒன்றிணைக்கிறது.
!!^a{formula} ஒன்று !!a{structure}~$\Struct{M}$ இல்,
மாறி ஒதுக்கீடு~$s$ சார்பாக நிறைவுறுத்தப்படலாம்; இதை
$\Sat{M}{!A}[s]$ என்று எழுதுகிறோம். இத்தொடர்பும்~$!A$ இன் கட்டமைப்பின்
மீது தொகுத்தறிதலால் வரையறுக்கப்படுகிறது; உதாரணமாக, தருக்கத் தொடுப்புக்குறிகளின்
மெய்மை அட்டவணைகளைப் பயன்படுத்தி $!A \land !B$ இன் நிறைவுறுத்தலை
$!A$ மற்றும் ~$!B$ ஆகியவற்றின் நிறைவுறுத்தல் (அல்லது நிறைவுறுத்தாமை) அடிப்படையில்
வரையறுக்கிறோம். மாறி ஒதுக்கீடு, !!{formula}~$!A$
!!a{sentence} ஆக, அதாவது கட்டற்ற மாறிகள் இல்லாததாக, இருந்தால் பொருத்தமற்றது
என்பது பின்னர் தெரியவரும்; எனவே !!{sentence}s !!{structure}s இல்
வெறுமனே நிறைவுறுத்தப்படுகின்றனவா இல்லையா என்று பேசலாம்.

% SEGMENT OLP-0645-S04
வாக்கியங்களுக்கான நிறைவுறுத்தல் தொடர்பு~$\Sat{M}{!A}$ அடிப்படையில்,
செல்லுபடியாகும் தன்மை, உட்கிடை, நிறைவுசெய்யத்தக்க தன்மை ஆகிய அடிப்படைப்
பொருண்மையியல் கருத்துகளை வரையறுக்கலாம். ஒரு வாக்கியம் செல்லுபடியாகும்,
$\Entails !A$, எனில் ஒவ்வொரு கட்டமைப்பும் அதை நிறைவுறுத்தும்.
!!{sentence}s கொண்ட ஒரு கணத்தால் அது உட்கிடையாகிறது,
$\Gamma \Entails !A$, எனில்~$\Gamma$ இல் உள்ள எல்லா !!{sentence}s ஐயும்
நிறைவுறுத்தும் ஒவ்வொரு !!{structure} உம்~$!A$ ஐயும் நிறைவுறுத்தும்.
மேலும், வாக்கியங்கள் கொண்ட ஒரு கணம் நிறைவுசெய்யத்தக்கது எனில், அதிலுள்ள
எல்லா !!{sentence}s ஐயும் ஒரே நேரத்தில் நிறைவுறுத்தும் !!{structure} ஒன்று
உள்ளது. !!{formula}s தொகுத்தறிதலால் வரையறுக்கப்பட்டுள்ளன; நிறைவுறுத்தலும்
!!{formula}s இன் கட்டமைப்பின் மீது தொகுத்தறிதலால் வரையறுக்கப்படுகிறது.
எனவே, நமது பொருண்மையியலின் பண்புகளை நிறுவவும் வரையறுக்கப்பட்ட
பொருண்மையியல் கருத்துகளைத் தொடர்புபடுத்தவும் தொகுத்தறிதலைப் பயன்படுத்தலாம்.

\end{document}
